\section{Finite-query normal form}

\subsection{The package}

We first state the normal form with one packed query variable.  Multiple
parameters are encoded before use; this loses no primitive-recursive
expressivity.

\begin{definition}[FQNF package]
An \(\FQNF\) package consists of fixed object-language data:
\begin{align*}
  &\Pre(q) &&\text{a precondition formula},\\
  &B(q) &&\text{a primitive-recursive certificate bound},\\
  &V(c,q) &&\text{a primitive-recursive Boolean verifier},\\
  &\Goal(c,q) &&\text{the conclusion extracted from an accepted code},
\end{align*}
together with object-iPRA derivations of:
\begin{align}
  &V(c,q)=0\;\vee\;V(c,q)=1,
      \tag{B}\label{eq:bool}\\
  &\Pre(q)\to\exists c\bigl(c<B(q)\wedge V(c,q)=1\bigr),
      \tag{P}\label{eq:prod}\\
  &V(c,q)=1\to\Goal(c,q).
      \tag{E}\label{eq:extract}
\end{align}
All free variables in these displays are universally available through the
usual open-axiom substitution mechanism.
\end{definition}

The verifier is not required to decide an arbitrary formula.  It is a fixed PR
program for one theorem family.  Its reflected meaning can be a rich bounded
formula or merely the local facts needed by the extractor.

\subsection{Selected certificate}

Let
\[
  \bad(c,q)=1\mono V(c,q),
\]
where \(\mono\) is truncated subtraction.  Under Booleanity,
\[
  \bad(c,q)=0 \quad\Longleftrightarrow\quad V(c,q)=1.
\]
Define \(\Sel(q)\) by bounded least search below \(B(q)\), returning \(B(q)\)
when no zero is found.  This is a fixed primitive-recursive term because the
verifier and bound are fixed.

\begin{theorem}[Finite-query compilation]\label{thm:fqnf-compile}
Every \(\FQNF\) package admits an object-iPRA derivation
\[
  \Pre(q)\to\Goal(\Sel(q),q).
\]
Moreover iPRA proves
\[
  \Pre(q)\to
  \bigl(\Sel(q)<B(q)\wedge V(\Sel(q),q)=1\bigr).
\]
\end{theorem}

\begin{proof}
Assume \(\Pre(q)\).  By the producer derivation
\eqref{eq:prod}, obtain a witness \(w\) satisfying
\(w<B(q)\) and \(V(w,q)=1\).  Booleanity and the defining arithmetic of
truncated subtraction give \(\bad(w,q)=0\).

Apply the checked bounded-search theorem to the body \(\bad(c,q)\), the bound
\(B(q)\), and witness \(w\).  It yields
\[
  \Sel(q)<B(q),\qquad \bad(\Sel(q),q)=0,
  \qquad \Sel(q)\le w.
\]
By Booleanity again, \(\bad(\Sel(q),q)=0\) implies
\(V(\Sel(q),q)=1\).  The extractor \eqref{eq:extract} gives
\(\Goal(\Sel(q),q)\).

The conclusion contains no occurrence of \(w\), so intuitionistic existential
elimination discharges it.  No excluded middle beyond the checker's proved
finite Boolean cases, and no choice principle, has been used.  Every step is a
constructor or checked arithmetic theorem of \(\Prf\,\PRAAx\).
\end{proof}

\begin{corollary}[Explicit witness form]
If a particular term \(w(q)\) is already proved valid and below the bound, then
iPRA proves the selected certificate valid and no later than \(w(q)\).
\end{corollary}

\begin{proof}
Use the same bounded-search theorem without existential elimination.  Its
minimality clause gives \(\Sel(q)\le w(q)\).
\end{proof}

\subsection{Conditional producer form}

Often the reusable verifier is elementary but existence depends on deep data.
Let \(D(q)\) be a boundary formula and replace the producer by
\[
  \Pre(q)\to D(q)\to
  \exists c<B(q)\;[V(c,q)=1].
\]
The same proof gives
\[
  \Pre(q)\to D(q)\to\Goal(\Sel(q),q).
\]
This is not a weaker engineering outcome.  It is the logically correct
factorization of, for example, a height-bound theorem from a finite point-list
checker.

\subsection{Why the theorem is genuinely constructive}

There are three possible sources of hidden classical strength, all absent here.

\begin{enumerate}
\item The witness is bounded and its predicate is checked by a fixed Boolean PR
term.
\item The selected code is a displayed PR term, not an application of an
abstract choice operator.
\item Existential elimination is used only to prove a conclusion independent
of the temporary witness.
\end{enumerate}

The theorem therefore scales uniformly with the query while remaining within
ordinary iPRA.
